Deepfakes threaten trust, safety, and platform integrity. Synthetic face\hyp swap videos and images, with generative models quickly becoming general purpose and easy to use, are increasingly weaponized for harassment, disinformation and misleading narratives, scams and frauds, intimate imagery without consent. Politically motivated deepfakes have been implicated in interfering with real\hyp world elections~\cite{chandra2025}; celebrity deepfakes are being weaponized in scams; synthetic identity fraud has cost institutions at least tens of millions of pounds. While conventional image manipulations disrupt the core tenets of authenticity, deepfakes produced by GANs or diffusion models (or increasingly complex reenactment systems) can ultimately reach photorealism both for human observers and for naive automated detectors.

Simultaneously, media consumption is becoming more heavily biased toward mobile devices---over 60\% of internet traffic is on phones, and most of TikTok \& Instagram's billions of users consume it through a mobile app, as does WhatsApp. Practical deepfake countermeasures must operate on underpowered phones/tablets, in a world of poor network access and strict privacy requirements. In many deployments, particularly privacy or security\hyp focused ones (as well as offline tools), no media leaves the device except the aggregate logs of detection signals. This on\hyp phone constraint precludes server\hyp side processing as a fallback on bad detections and forces models to be simultaneously accurate, compact, and fast.

Mobile platforms therefore inhabit a fundamental tension between strong detection accuracy and a user\hyp friendly experience. The model must have a high recall (low false negative rate), ensuring harmful content does not evade detection by the model, but at the same time it must also be small enough to not break app size budgets and fast enough to avoid noticeable latency and battery drain. Under distribution shift---where it encounters manipulation techniques, compression pipelines, and content types not seen during training---detectors should still be robust without reverting to expensive fallback models for all inputs. We focus on this on\hyp device deployment setting, and the three primary challenges it entails: (i)~cross\hyp dataset generalization in the wild, (ii)~low latency and small model footprint on commodity hardware, and (iii)~control of false negatives (missed fakes) under operational compute budgets.

\paragraph{Deepfake detection research covers a vast space.}
Deepfake detection research covers a vast space including spatial, frequency, and temporal cues, multimodality, dataset curation, and evaluation protocols. Recent comprehensive surveys~\cite{verdoliva2020deepfakes,pei2024deepfakesurvey,chen2024deepfakes} document two longstanding gaps which motivate our work: (i)~detectors that perform well on curated benchmarks often severely under\hyp perform under distribution shift, and (ii)~the majority of evaluation assumes server class hardware as opposed to mobile devices. These gaps imply an opportunity for system designs that consider cross\hyp dataset generalization and mobile constraints as first\hyp class goals rather than an afterthought.

\paragraph{Deepfake Detection as a Problem Space.}
Deepfake detection is then a binary classification problem of whether given an image or frame contains a real or synthetically manipulated face. While this sounds simple, the class member properties are produced by diverse manipulation techniques (face swapping, face reenactment, attribute editing, full synthesis) while the environmental properties (compression artifacts, lighting, camera sensors, post\hyp processing) also vary widely. Early works used hand\hyp crafted forensic features like head pose inconsistencies~\cite{yang2019}, blinks of the eyes, JPEG compression artifacts~\cite{cozzolino2019}, while more modern methods let deep learning automatically extract hierarchical representations directly from pixels or Fourier/frequency spectra producing state of the art accuracies on in\hyp distribution benchmarks~\cite{rossler2019faceforensicspp,dolhansky2020deepfake,li2020celebdf}.

Cross\hyp dataset generalization is a distant goal: models trained on FaceForensics++ or CelebDF suffer AUC drops of 20\hyp 30\% evaluated on in the wild datasets such as WildDeepfake~\cite{zi2021wilddeepfake} or Deepfake\hyp Eval\hyp 2024~\cite{chandra2025}. This generalization gap is primarily the result of dataset bias (e.g.\ training set would overrepresent certain generators/compression settings), domain shift (e.g.\ social media would apply proprietary compression pipelines) and the rapid evolution of synthesis techniques. Bridging these gaps requires training protocols on multiple datasets exposing models to diverse artifacts, as well as trustworthy cross\hyp dataset protocols that allow out\hyp of\hyp distribution/OOD generalization from studying the robustness of such methods beyond just in\hyp distribution accuracy.

\paragraph{Mobile Deployment Constraints.}
Most deepfake detection research operates under the assumption of unlimited computation resources---models are trained and evaluated on server GPUs with batch sizes and latencies that are untenable for on\hyp device inference. Many real\hyp world use cases (live video verification, client\hyp side content moderation, edge forensics, etc.) must be deployed on commodity smartphones or embedded devices with limited memory, power budgets, and thermal envelopes. Mobile deployment comes with additional wrinkles: (i)~model size must reside within common app bundle size constraints (50--100\,MB), (ii)~per\hyp frame latency should be amenable to real\hyp time or near\hyp real time inference (e.g., $<$\,200\,ms), and (iii)~energy draw should maintain battery life over hours\hyp long usage.

Several strands of prior work are relevant to this goal. Deepfake\hyp specific detectors explore spatial, frequency, and physiological cues~\cite{pei2024deepfakesurvey,durall2020watch,qian2020thinkingfrequency}, while lightweight vision models such as MobileNet and EfficientNet families target efficient inference on ARM\hyp class hardware~\cite{mobilenetv4,efficientnetv2,zhang2018shufflenet,ma2018shufflenetv2}. Separately, early\hyp exit networks and cascaded systems~\cite{teerapittayanon2017branchynet,huang2018msdnet,kaya2019shallowdeep} show how to trade compute for accuracy by routing easy samples to shallow branches. However, existing deepfake work rarely combines these ingredients into an end\hyp to\hyp end, auditable system that is (i) trained across multiple datasets, (ii) tuned with explicit risk and cost metrics, and (iii) packaged as a reproducible mobile deployment template. Most studies either evaluate a single heavy detector on server\hyp class GPUs or treat mobile optimization as a brief post\hyp processing step.

In this work we adopt a system\hyp level perspective and design a cascaded detector tailored for on\hyp device deployment. Our goal is to expose explicit knobs over recall, escalation rate, and compute while keeping the training and evaluation pipeline reproducible from raw manifests to mobile artifacts. Against this backdrop, we ask:
\begin{itemize}[leftmargin=*]
  \item \emph{How can we design a deepfake detector that provides high recall under distribution shift while remaining deployable on resource\hyp constrained mobile devices?}
  \item \emph{Can a cascaded architecture with explicit thresholds offer interpretable trade\hyp offs between compute and risk, without sacrificing reproducibility?}
  \item \emph{What minimal set of scripts and artifacts is needed so that practitioners can audit, reproduce, and extend the full pipeline?}
  \item \emph{How should we evaluate and report robustness---for example under compression, noise, and blur---and cross\hyp dataset generalization in a way that directly informs deployment decisions?}
\end{itemize}
Existing works on mobile vision and early exits provide building blocks (lightweight backbones, quantization, cascades), but rarely integrate them into an end\hyp to\hyp end, auditable system tailored to deepfake detection with multi\hyp dataset training and out\hyp of\hyp distribution (OOD) evaluation.

\paragraph{Contributions.} We design \textbf{MobileDeepfake}, a cascaded detection system engineered for mobile deployment and rigorous evaluation. Our contributions are:
\begin{itemize}[leftmargin=*]
  \item \textbf{A cost\hyp aware, two\hyp stage cascade with explicit risk controls.} A lightweight MobileNetV4 Stage~1 acts as a fast filter; ambiguous samples are escalated to an EfficientNetV2 expert (Stage~2), with an optional Stage~3 LightGBM meta\hyp model. We define Stage~1 leakage (the fraction of fakes incorrectly accepted at Stage~1) and Stage~2 escalation rate as primary system metrics, and use a simple two\hyp threshold routing scheme coupled with a cost\hyp sensitive grid search (formalised in Section~\ref{sec:problem_formulation}) to expose interpretable trade\hyp offs between false negative rate (FNR) and compute.
  \item \textbf{Multi\hyp dataset training and cross\hyp dataset evaluation, including Deepfake\hyp Eval\hyp 2024.} We train on four academic datasets and hold out Deepfake\hyp Eval\hyp 2024 for OOD validation and test, quantifying performance gaps and highlighting practical failure modes under in\hyp the\hyp wild distribution shift. Our experiments report both single\hyp model baselines and cascaded results, with calibrated vs.\ uncalibrated comparisons when transferring thresholds across datasets.
  \item \textbf{An end\hyp to\hyp end, scriptable pipeline with hard\hyp example mining, calibration, and robustness analysis.} The system includes data preprocessing, manifest generation, hard\hyp sample mining, expert training, meta\hyp model fitting, and threshold tuning, with scripts that generate the tables and figures used in this paper. Robustness sweeps over common corruptions (JPEG compression, noise, blur, brightness) and a Deepfake\hyp Eval\hyp 2024 evaluation stage are integrated into the same pipeline.
  \item \textbf{A mobile\hyp oriented export and integration template.} We provide a mobile optimization stage that applies quantization and knowledge distillation to shrink models, and a deployment stage that exports compact TorchScript artifacts (and ONNX bundles when needed) together with thresholds/temperature in a metadata file suitable for Android integration. While our focus is on methodology rather than a new backbone, the implementation illustrates how to turn deepfake research pipelines into deployable, auditable systems.
\end{itemize}

Concretely, MobileDeepfake follows a train\,\textrightarrow\,tune\,\textrightarrow\,deploy workflow across six stages. The core online inference in our final system is a two\hyp stage cascade: Stage~1 trains a lightweight MobileNetV4 filter for fast screening; Stage~2 builds an EfficientNetV2 expert and produces per\hyp sample logits/features. Stage~3 optionally fits a LightGBM meta\hyp model on K\hyp fold meta\hyp datasets to study fusion but is not part of the default deployment. Stage~4 grid\hyp searches routing thresholds to minimize false negatives under an escalation\hyp rate budget and logs Stage~1 leakage alongside Stage~2 usage. Stage~5 conducts cross\hyp dataset and robustness evaluation (including Deepfake\hyp Eval\hyp 2024), and Stage~6 exports quantized artifacts for on\hyp device integration.

We focus on deployability rather than proposing a new backbone: the contribution is a practical, auditable methodology that couples a two\hyp stage cascade with explicit threshold controls, calibration, robustness analysis, and a mobile export path. We do explore stronger hybrid CNN\hyp Transformer experts such as GenConViT in our Stage~2 ensemble, but find that they offer limited and unstable gains relative to EfficientNetV2\hyp B3 under our constraints, so we relegate them to an appendix and base the default cascade on MobileNetV4 + EfficientNetV2. Subsequent sections detail the datasets and task, the cascaded method and optimization objectives, the experimental setup, and empirical results on in\hyp distribution and out\hyp of\hyp distribution data, followed by a discussion of limitations and deployment considerations.
